\chapter{一束光到一圈扫描：CPU 激光雷达}\label{ch:lidar}
\begin{notebox}
\textbf{目标}\quad 用可手算的几何求交构建激光雷达，理解遮挡、角分辨率与测距误差。\\
\textbf{先修}\quad 二维向量、圆方程，以及前四章的坐标、世界与机器人模型。
\end{notebox}

\section{先让一束光碰到障碍}
雷达安装在圆盘中心，发出一条没有宽度的射线。机器人位置与 yaw 决定射线的起点和方向。
雷达不检测机器人自身，否则从圆心发出的光总会先碰到自己的外壳。机器人半径仍参与运动碰撞。

射线用起点 o、单位方向 d 和非负距离 t 表示：
\begin{equation}
p(t)=o+td,\qquad \|d\|=1,\quad t\geq0.
\end{equation}
只有方向已归一化时，t 才直接表示米。测距求的是最近表面，不是障碍中心的距离。
沿同一方向有多个交点时，较近表面会挡住较远表面。

\begin{center}
\begin{tikzpicture}[x=1cm,y=1cm]
\draw[fill=tealblue!15,draw=tealblue,thick] (3,0) circle (1);
\draw[navy,very thick] (5,-1.3)--(5,1.3);
\draw[-{Stealth},navy,thick] (0,0)--(2,0);
\draw[dashed,muted] (2,0)--(5,0);
\fill[navy] (0,0) circle (.05) node[below] {$o$};
\fill[tealblue] (2,0) circle (.06) node[above left] {最近命中};
\fill[navy] (3,0) circle (.04) node[above] {$c$};
\node[above] at (5,1.3) {被遮挡的墙};
\node[below] at (1,-.35) {$t=2$ m};
\end{tikzpicture}
\end{center}

重新构建工作空间并加载 install，先运行不依赖 ROS 回调的几何演示：
\begin{lstlisting}[language=bash,numbers=none]
ros2 run motion2d raycast_demo
\end{lstlisting}
三个距离应依次为 2、4、2 m。前两个分别是圆和墙的单独求交；最后一个在世界中取最近值。
代码位于 \path{src/examples/raycast_demo.cpp}，这里的墙和圆都可以手算。

\clearpage
\section{射线与圆：别忘了第二个根}
把射线代入圆心 c、半径 R 的圆方程，令 $q=c-o$、$h=q^\top d$：
\begin{equation}
\|o+td-c\|^2=R^2
\ \Longrightarrow\ t^2-2ht+\|q\|^2-R^2=0.
\end{equation}
二维叉积 $a\times b=a_xb_y-a_yb_x$ 是一个标量。
单位方向下，$|d\times q|$ 是圆心到射线所在直线的距离，因此
\begin{equation}
D=R^2-(d\times q)^2,\qquad t_\pm=h\pm\sqrt D.
\end{equation}
$D<0$ 时没有交点；$D=0$ 时相切，也算命中。先检查较小的根 $t_-$ 是否非负，
否则再检查 $t_+$。两个根都负，说明圆在射线背后。圆内起点的小根为负，大根是出口；
这有助于测试函数，但课程中的机器人应始终位于自由空间。

\par\noindent\begin{minipage}{\linewidth}
\lstinputlisting[language=C++,linerange={ray_circle_begin-ray_circle_end},
  includerangemarker=false]{../ros2_ws/src/motion2d/src/sim/raycast.cpp}
\end{minipage}\par

没有命中返回正无穷，它与后续所有距离取最小值时自然不起作用。此处先计算真实几何，
量程与测量噪声留给扫描层处理；不要在这个函数里混入 ROS 消息。

\textbf{代码练习 ch05-1：}在 \path{exercises/ch05/starter/forward_root.hpp}
补完根的选择。从仓库根目录运行：
\begin{lstlisting}[language=bash,numbers=none]
g++ -std=c++17 -Iexercises/ch05/starter \
  exercises/ch05/check_root.cpp -o /tmp/ch05_root
/tmp/ch05_root
\end{lstlisting}
先手算圆心 (3,0)、半径 1、起点 (0,0)、方向 (1,0) 的两个根；再把起点放到圆心。
这两个例子分别检查最近回波和第二个根。

\clearpage
\section{射线与线段：平行不一定没交点}
令线段端点为 a、b，$e=b-a$。联立 $o+td=a+ue$，分别取叉积可得
\begin{equation}
t=\frac{(a-o)\times e}{d\times e},\qquad
u=\frac{(a-o)\times d}{d\times e}.
\end{equation}
只有 $t\geq0$ 且 $0\leq u\leq1$ 才命中线段，而不只是命中它所在的直线。
当分母接近零时不能继续做除法：平行但不共线则无交点；共线时将端点投影到射线上，
检查投影区间是否与 $[0,+\infty)$ 相交，取重叠部分的最近点。

\par\noindent\begin{minipage}{\linewidth}
\lstinputlisting[language=C++,linerange={ray_segment_begin-ray_segment_end},
  includerangemarker=false]{../ros2_ws/src/motion2d/src/sim/raycast.cpp}
\end{minipage}\par

世界求交 \code{raycastWorld} 遍历所有圆、所有多边形边，以及矩形边界的四条边，
对得到的非负距离取最小值。这样自然处理遮挡，不依赖障碍物存储顺序。
算法复杂度随表面数线性增长；先保留这个清楚的 CPU 基准，后续再讨论加速。

\textbf{验收与实验：}运行 \code{colcon test --packages-select motion2d}。
几何测试覆盖圆相切、背向、线段端点、平行与共线、矩形角点、遮挡和刚体变换不变性。
把演示的圆移到墙后，预测最近命中；反转射线方向，解释距离为何改变。

\begin{notebox}
这里的浮点容差面向课程的米制小场景。它用于避免接近平行时除以极小数，
不是任意坐标尺度下的精确几何谓词。无效几何应在世界生成边界拒绝。
\end{notebox}
